\section{Introduction}

Accurate indoor positioning is an important enabling technology for location-aware Internet of Things (IoT) services, asset tracking, healthcare monitoring, navigation, industrial automation, and intelligent buildings. Although Global Navigation Satellite Systems provide reliable positioning in outdoor environments, their performance deteriorates significantly indoors because satellite signals are attenuated by walls, floors, and other structural obstacles. Consequently, a wide variety of alternative indoor positioning technologies have been investigated, including ultra-wideband, Bluetooth Low Energy, inertial sensing, visual positioning, and Wi-Fi-based localization \cite{qiao2024survey,dai2023survey}. Among these alternatives, Wi-Fi is particularly attractive because positioning can often reuse communication infrastructure that is already deployed in residential, commercial, and institutional buildings.

Traditional Wi-Fi fingerprinting systems primarily rely on Received Signal Strength (RSS). A radio map is constructed during an offline phase by collecting RSS measurements at known reference points (RPs), while during online positioning the observed fingerprint is matched or mapped to an estimated user location. Despite its widespread availability, RSS is highly sensitive to multipath propagation, device heterogeneity, human movement, environmental changes, and temporal signal fluctuations \cite{shang2022overview,feng2022deep}. These effects can significantly degrade positioning accuracy and reduce the transferability of models across different times and environments.

The introduction of Fine Timing Measurement (FTM) in IEEE 802.11mc enabled commodity devices to obtain Wi-Fi Round-Trip Time (RTT) measurements. Unlike RSS, which provides an indirect indication of distance through received signal power, RTT provides a time-of-flight-related ranging observable. Recent research has therefore increasingly investigated RTT-based positioning and the fusion of RTT with RSS information \cite{martinescalona2023ftm,eberechukwu2023fusion,feng2023rttfingerprinting,dong2023rtt_rss}. RTT is nevertheless affected by non-line-of-sight (NLOS) conditions, multipath propagation, access-point geometry, and measurement availability, motivating dedicated LOS/NLOS compensation and robust learning methods \cite{cao2024los,lyu2025visibility}.

Recent studies demonstrate that combining RTT and RSS can provide more robust positioning information than relying on either modality independently. Feng et al. \cite{feng2026robust} combined RTT and RSS fingerprints with Random Forest regression and conformal prediction, obtaining sub-meter positioning performance across multiple indoor environments while providing uncertainty regions around the estimated positions. Rana and Dulal \cite{rana2025gtcn} subsequently investigated a graph--temporal neural architecture that jointly models AP geometry and hybrid RTT--RSS measurements. At the same time, another important research direction has focused on reducing positioning complexity and Wi-Fi RTT communication overhead. Gonzalez Diaz et al. \cite{gonzalezdiaz2026scalable,gonzalezdiaz2024stdclustering} showed that AP selection, clustering, and lightweight proximity-based processing can significantly reduce computational and ranging requirements, which is particularly important for resource-constrained IoT devices.

These studies highlight an important trade-off. Increasing model complexity or exploiting more measurements may improve positioning performance, but it can also increase training cost, inference complexity, model size, and network ranging overhead. Conversely, aggressively reducing the measurement or feature space can improve efficiency at the cost of localization accuracy. Recent studies on AP selection and feature selection confirm that not every available Wi-Fi measurement contributes equally to localization \cite{wang2024gnnap,ayinla2026reloc}. Redundant, correlated, unstable, or weakly informative AP measurements may increase model complexity without providing proportional improvements in localization performance. The same redundancy matters not only for which measurements are collected, but also for how a learner uses them.

Random Forest regression provides an attractive compromise for fingerprint-based indoor localization because it offers nonlinear modeling capability, limited preprocessing requirements, robustness to heterogeneous feature distributions, low online complexity, and training times of seconds. In recent hybrid RTT--RSS work, however, the forest is typically used with every feature eligible at every split~\cite{feng2026robust}. This setting, the default of the \texttt{scikit-learn} regressor, is bagging of regression trees rather than a random forest in Breiman's sense~\cite{breiman2001random}: the per-split feature subsampling that decorrelates the trees is switched off. Hybrid fingerprints are a setting in which this matters, because the RTT and RSS of one AP, and the measurements of neighboring APs, carry overlapping spatial information, while NLOS propagation and visibility changes make individual features intermittently unreliable. Unrestricted split selection then lets the same dominant measurements shape most trees.

This work therefore studies the number of candidate features per split, $m_{\mathrm{try}}=\rho p$, as the central design variable of hybrid RTT--RSS forests. Feature subsampling itself is not new~\cite{ho1998random,breiman2001random}, and it is known to act as a regularizer whose optimal strength depends on the signal-to-noise ratio of the problem~\cite{mentch2020randomization,probst2019hyperparameters}. What is not known for Wi-Fi fingerprinting is (i) how much it matters relative to the full-feature setting used in practice, (ii) which ratio should be used and how it can be chosen without access to test locations, and (iii) how the answer depends on the density of the surveyed radio map, which is the main cost driver of fingerprinting deployments. We use a fixed default $\rho=0.7$, frozen before any external evaluation, and analyze it against a full ratio sweep, alternative selection rules, and alternative tree ensembles.

A second issue concerns evaluation. Fingerprinting datasets contain many repeated scans at each physical reference point (RP). The official splits of the benchmark used here hold out complete RPs, but sample-level protocols, which place scans of the same RP in both training and validation, remain common, including in the out-of-bag error and cross-validation routines used to tune forests. Such protocols measure recognition of already-surveyed locations rather than generalization to new ones~\cite{roberts2017crossvalidation,wang2024domaininvariant,long2025multibuilding,narasimman2024dumbloc}. We therefore group all scans by their physical RP and report sample-level and RP-disjoint results separately; as shown below, this distinction also decides how the subspace ratio can be chosen.

The same correlation affects uncertainty quantification. Conformal prediction has been introduced for hybrid RTT--RSS positioning~\cite{feng2026robust} and for RSS fingerprinting~\cite{wang2023maploc,zhang2024rloc,ma2024navigating,nipa2025uncertainty,guo2026uegloc}, but its finite-sample guarantee assumes exchangeable calibration samples, whereas the $\sim$120 scans of one RP are nearly duplicates. We show that this matters when few RPs are available for calibration and compare scan-level calibration with a valid group construction and a simple RP-count correction.

A third issue is external validity. Improvements obtained on one benchmark may reflect benchmark-specific tuning, so the configuration is also evaluated, without retuning, on the independent EETAC auditorium dataset~\cite{martinescalona2023ftm}, which has also been used for scalable RTT positioning~\cite{gonzalezdiaz2026scalable}. In doing so, we found that the two EETAC acquisition days are recorded in mirrored coordinate frames; we detect and correct this from the data before any pooled or cross-day evaluation.

The scope of this work should be stated clearly: it concerns inexpensive tree-based models, and it does not claim state-of-the-art accuracy. On the Building environment, the graph--temporal network GTCN~\cite{rana2025gtcn} reports a 2D RMSE of $0.660$~m, whereas the proposed forest reaches $0.732$~m. The two approaches occupy different points of the design space: GTCN requires about 38~min of training for 1500 epochs, whereas the forests train in about 1--15~s on a laptop CPU, need no hyperparameter search beyond the subspace ratio, and can be retrained immediately when the radio map is resurveyed. For deployments in which the radio map changes, or in which training must run on modest hardware, this cost difference matters as much as the accuracy gap, and the question studied here, how to configure and evaluate such inexpensive models correctly, remains practically relevant.


\subsection{Main Contributions}

Within this scope, the main contributions of this work are:

\begin{itemize}

	\item \textbf{Feature subsampling versus radio-map density.}
	Across three environments and four survey densities, full-feature bagging ($\rho=1$) is never the best ratio, and the optimal ratio decreases as the radio map becomes sparser in all three environments (Building: 0.5 $\rightarrow$ 0.2; Apartment: 0.85 $\rightarrow$ 0.35; Office: 0.7 $\rightarrow$ 0.35), while the gain of subsampling over bagging tends to increase (up to 25.5\%). Extremely randomized trees, which already randomize split thresholds, gain at most 0.8\% from subsampling. This explains why validation protocols that train on fewer RPs prefer smaller ratios than the full-density test.

	\item \textbf{How to choose the ratio.}
	We compare out-of-bag error, sample-level CV, RP-grouped CV, leave-one-environment-out selection, and fixed defaults. Data-driven rules are biased toward small ratios (worst-case test regret 6.0--22.5\%), whereas the fixed default $\rho=0.7$ stays within 4.1\% of the best ratio in every environment (bagging: 19.7\%).

	\item \textbf{Accuracy over the reproduced baseline and literature ensembles.}
	The reproduced full-feature RF matches the published values; $\rho=0.7$ reduces its paper-compatible error by 14.3\%, 7.1\%, and 5.1\% on the official test sets, with RP-clustered bootstrap confidence intervals, and is better in all ten seeds in every environment. Breiman's $p/3$ default, extremely randomized trees, Rotation Forest, gradient-boosted trees, and WKNN are evaluated under the same protocols.

	\item \textbf{Corrected independent validation.}
	On the EETAC dataset, after aligning the mirrored acquisition-day frames, the frozen configuration reduces 2D RMSE by 11--14\% under sample-level and RP-disjoint CV and is better than the baseline in 74 of 80 seed--protocol comparisons; cross-day errors are 1.2--1.7~m instead of more than 5~m without alignment, and one cross-day direction is a tie.

	\item \textbf{Compact deployment.}
	Because subsampling buys accuracy, 50--100-tree forests remain more accurate than the 300-tree full-feature model while being 2.5--5.5 times smaller and faster at inference.

	\item \textbf{Calibration with clustered scans.}
	With 21 held-out RPs, scan-level split-conformal calibration falls below the nominal coverage on average and in 65--75\% of random calibration splits. Group calibration with one scan per RP~\cite{Dunn02102023} attains the nominal mean coverage with 22--30\% larger regions and needs at least $1/\alpha-1$ calibration RPs, which gives a survey-planning rule; a simpler RP-count correction is conservative. Covariance-aware elliptical regions are analyzed as an anisotropic alternative to circles.

\end{itemize}

All experiments, including the extension analyses, are produced from the public data by released scripts, and a notebook reproduces every table and figure (see Code and Data Availability).


\subsection{Paper Organization}

Section~\ref{sec:related_work} reviews RSS and RTT fingerprinting, hybrid RTT--RSS localization, tree-ensemble randomization, and uncertainty-aware positioning. Section~\ref{sec:system_model} defines the system model. Section~\ref{sec:proposed_method} presents the subspace forest, the choice of the subspace ratio, RP-aware validation, and the uncertainty constructions. Section~\ref{sec:experimental_setup} describes the datasets, preprocessing, protocols, and metrics. Section~\ref{sec:results} reports the results, Section~\ref{sec:discussion} discusses them, Section~\ref{sec:limitations} lists the limitations, and Section~\ref{sec:conclusion} concludes the paper.
